\documentclass[a4paper, 11pt]{article}
\usepackage{inputenc}
\usepackage{amsmath}
\usepackage{graphicx}
\usepackage[T1]{fontenc}
\usepackage{babel}
\usepackage{hyperref}    %pour indexer la table des matières 
\hypersetup{pdfborder=1 1 1}
\usepackage{float} 
\usepackage{url}
\usepackage{caption}
\usepackage{multirow}
\usepackage{multicol}
\usepackage{listings}
\usepackage{subcaption}
\usepackage{empheq}
\usepackage{fancyhdr}
\usepackage{booktabs}
\usepackage{amssymb}
\usepackage{xcolor}
\usepackage{setspace}
\usepackage{enumitem}
\setlength{\hoffset}{-18pt}         
\setlength{\oddsidemargin}{0pt} % Marge gauche sur pages impaires
\setlength{\evensidemargin}{9pt} % Marge gauche sur pages paires
\setlength{\marginparwidth}{54pt} % Largeur de note dans la marge
\setlength{\textwidth}{481pt} % Largeur de la zone de texte (17cm)
\setlength{\voffset}{-18pt} % Bon pour DOS
\setlength{\marginparsep}{7pt} % Séparation de la marge
\setlength{\topmargin}{0pt} % Pas de marge en haut
\setlength{\headheight}{13pt} % Haut de page
\setlength{\headsep}{10pt} % Entre le haut de page et le texte
\setlength{\footskip}{27pt} % Bas de page + séparation
\setlength{\textheight}{708pt} % Hauteur de la zone de texte (25cm)
%Traits en haut et en bas des pages
\usepackage{fancyhdr}
\pagestyle{fancy}
\fancyhead[L]{}
\fancyhead[R]{}
\renewcommand\footrulewidth{1pt}
\renewcommand\footrulewidth{1pt}
\fancyfoot[L]{\tiny Project Report}
\usepackage{lastpage}
\fancyfoot[R]{\tiny January 2024}
\fancyfoot[C]{\textbf{Page \thepage/\pageref{LastPage}}}
\lstnewenvironment{mathematicacode}[1][]
{
	\lstset{
		language=Mathematica,
		basicstyle=\small\ttfamily,
		numbers=left,
		numberstyle=\tiny,
		numbersep=5pt,
		frame=single,
		frameround=tttt,
		framexleftmargin=5pt,
		#1
	}
}
{}

\begin{document} 
	
	%----------------------------------------------------------------------------------------
	%	TITLE PAGE
	%----------------------------------------------------------------------------------------
	
	\begin{titlepage} % Suppresses displaying the page number on the title page and the subsequent page counts as page 1
		\newcommand{\HRule}{\rule{\linewidth}{0.5mm}} % Defines a new command for horizontal lines, change thickness here
		
		\centering % Centre everything on the page
		
		%------------------------------------------------
		%	Headings
		%------------------------------------------------
		
		\textsc{\LARGE ENSEIRB-MATMECA}\\[1cm] % Main heading such as the name of your university/college
		
		
		\begin{figure}[H]
			\begin{center}
				\includegraphics[scale=4]{logo ecole.jpg}
			\end{center}
		\end{figure}
		
		%\textsc{\large Minor Heading}\\[0.5cm] % Minor heading such as course title
		
		%------------------------------------------------
		%	Title
		%------------------------------------------------
		
		\vspace{1cm}
		
		\HRule\\[0.4cm]
		
		{\huge\bfseries Solving the linear Boltzmann equation using Monte Carlo methods }\\[0.4cm] % Title of your document
		
		\HRule\\[1.5cm]
		
		{\huge\bfseries }
		
		\bigskip
		\bigskip
		
		\centering
		\Large{Project lead by :} \\
		\Large{C. Aumonier, A. Boucher, G. Doyen, K. El Maddah, G. Rodiere}
		
		%------------------------------------------------
		%	Date
		%------------------------------------------------
		
		\vfill\vfill\vfill % Position the date 3/4 down the remaining page
		{\Large  Project supervised by :}\\
		\Large{G. Poëtte}
		
		\vspace{0,5cm}
		
		{\large January 2024} % Date, change the \today to a set date if you want to be precise
		
		%------------------------------------------------
		%	Logo
		%------------------------------------------------
		
		%\vfill\vfill
		%\includegraphics[width=0.2\textwidth]{placeholder.jpg}\\[1cm] % Include a department/university logo - this will require the graphicx package
		
		%----------------------------------------------------------------------------------------
		
		\vfill % Push the date up 1/4 of the remaining page
		
	\end{titlepage}
	
\tableofcontents

\newpage
	
\section{Introduction}

**Gabriel

\paragraph{Note} Here is a brief summary of the various existences and uniqueness of solutions based on different cases of equations to be solved. The primary source used is the transport and diffusion document.

\section{Transport Equation}

\paragraph{}
In this section, we will examine the solutions of the transport equation, a particular case of the Boltzmann equation, defined below:



Let's $v \in \mathbb{R}^N$

\begin{equation} \label{transport}
	\frac{\partial f}{\partial t}(t,x)+v \cdot \nabla_x f(t,x) + a(t,x)f(t,x) = S(t,x), \quad x \in \mathbb{R}^N, \quad t>0.
\end{equation}


$f$ is the unknown function, $a$ is the given damping or amplification rate, and $S$ is the given source term.

\subsection{Cauchy Problem}
\subsubsection{Special Case without Absorption or Source Term}

Let's consider the following Cauchy problem:

\[
\begin{cases}
	\frac{\partial f}{\partial t}(t,x)+v \cdot \nabla_x f(t,x)=0, \quad x \in \mathbb{R}^N, \quad t>0.\\
	f(0,x)) = f^{in}(x)
\end{cases}
\]

If $f^{in} \in \mathcal{C}^1(\mathbb{R}^n)$, then using the method of characteristics, we can easily demonstrate that the Cauchy problem has a unique solution $f \in \mathcal{C}^1(\mathbb{R}_+ \times \mathbb{R}^n)$. This solution is given as follows:





$\boxed{f(t,x)=f^{in}(x-tv) \quad \forall (t,x) \in \mathbb{R}_+ \times \mathbb{R}^n }$




\subsubsection{With Absorption and Source Term}

This time, we consider damping and the source term, assuming they belong to $\mathcal{C}^1(\mathbb{R}_+ \times \mathbb{R}^n)$. The considered Cauchy problem is as follows:


\[
\begin{cases}
	\frac{\partial f}{\partial t}(t,x)+v \cdot \nabla_x f(t,x) + a(t,x)f(t,x) = S(t,x), \quad x \in \mathbb{R}^N, \quad t>0.\\
	f(0,x)) = f^{in}(x)
\end{cases}
\]

In this case, if $f^{in} \in \mathcal{C}^1(\mathbb{R}^n)$, then using the method of characteristics, we can still demonstrate that the Cauchy problem has a unique solution $f \in \mathcal{C}^1(\mathbb{R}_+ \times \mathbb{R}^n)$. This solution is given as follows:\\
$\boxed{f(t,x)=f^{in}(x-tv) e^{-\int_0 ^t a(\tau,x+(t-\tau)v)d\tau} + \int_0 ^t  e^{-\int_0 ^t a(\tau,x+(t-\tau)v)d\tau} S(s,x+(s-t)v)ds \quad \forall (t,x) \in \mathbb{R}_+ \times \mathbb{R}^n }$

\subsection{Boundary Value Problems}

In the previous section, we presented results applicable when considering complete freedom for the spatial variable in $\mathbb{R}^n$. However, for a significant portion of physical problems, this may not be the case. That is why we are more interested in a problem of the boundary value type. In this section, we will consider the set $\Omega$, which is an open set with a $\mathcal{C}^1$ boundary in $\mathbb{R}^n$. Additionally, we take $v \in \mathbb{R}^n\backslash \{0\}$.

\subsubsection{Special Case without Absorption or Source Term}

\paragraph{}
If, moreover, $\Omega$ is convex, and let $f_b^- \in \mathcal{C}^1(\mathbb{R}_+ \times \partial \Omega^-)$ and $f^{in} \in \mathcal{C}^1(\overline{\Omega})$, where $\partial \Omega^- = \{x\in\partial \Omega~ | ~ v\cdot n_x<0\}$, then the problem:


\[
\begin{cases}
	\frac{\partial f}{\partial t}(t,x)+v \cdot \nabla_x f(t,x)=0, \quad x \in \mathbb{R}^N, \quad t>0.\\
	\restriction{f}{\partial \Omega}= f_b^-\\
	\restriction{f}{t=0}= f^{in}(x)
\end{cases}
\]

has a unique solution $f \in \mathcal{C}^1(\mathbb{R}_+ \times \Omega)$ and $f \in \mathcal{C}^0(\mathbb{R}_+ \times \overline{\Omega})$ if and only if $\forall y \in \partial \Omega^-,$

\[ f_b^-(0,y) = f^{in}(y), \quad \frac{\partial f_b^-}{\partial t}(t,y)+v \cdot \nabla f^{in}(y)=0. \]

\paragraph{}

The expression of the solution is then given by: 
\[
\boxed{
	f(t,x)= 
	\begin{cases}
		f^{in}(x-tv) \quad si~ t \leq \tau_x ,\\
		f_b^-(t-\tau_x,x-\tau_xv) \quad si~ t > \tau_x ,
	\end{cases}
}
\]

With $\tau_x= \inf \{t \geq 0~ | ~ x-tv \notin \overline{\Omega}\}$

\subsubsection{With Absorption and Source Term}

\paragraph{}
We still consider $\Omega$ to be convex, with $f_b^- \in \mathcal{C}^1(\mathbb{R}_+ \times \partial \Omega^-)$, $f^{in} \in \mathcal{C}^1(\overline{\Omega})$, $a$, and $S$ $\in \mathcal{C}^1(\mathbb{R}_+ \times \overline{\Omega})$. If we assume that $\forall y \in \partial \Omega^-$,

\[
\begin{aligned}
	f_b^-(0,y) &= f^{in}(y), \\
	\frac{\partial f_b^-}{\partial t}(0,y) + v \cdot \nabla_x f^{in}(y) + a(0,y)f^{in}(y)&=S(0,y)
\end{aligned}
\]



In this case, the problem becomes:

\[
\begin{cases}
	\frac{\partial f}{\partial t}(t,x)+v \cdot \nabla_x f(t,x)+ a(t,x)f(t,x)=S(t,x), \quad x \in \mathbb{R}^N, \quad t>0.\\
	\restriction{f}{\partial \Omega}= f_b^-\\
	\restriction{f}{t=0}= f^{in}(x)
\end{cases}
\]

has a unique solution $f \in \mathcal{C}^1(\mathbb{R}_+ \times \Omega)$ and $f \in \mathcal{C}^0(\mathbb{R}_+ \times \overline{\Omega})$. The expression for $f$ is then:


\[
\boxed{
	\begin{aligned}
		f(t,x) &= \mathbf{1}_{t \leq \tau_x} f^{in} (x-tv) \exp\left(-\int_0^t a(s,x+(s-t)v)ds\right) \\
		&\quad + \mathbf{1}_{t > \tau_x} f_b^- (t-\tau_x,x-\tau_x v) \exp\left(-\int_{t-\tau_x}^t a(s,x+(s-t)v)ds\right) \\
		&\quad + \int_{t-\tau_x}^t \exp\left(-\int_{s}^{t}a(\tau,x+(\tau-t)v)d\tau\right) S(s,x+(s-t)v)ds
\end{aligned}}
\]


\subsection{Generalized Solutions}

In this section, $\Omega$ is not necessarily convex. This implies that the continuity of potential solutions is not guaranteed. However, in practice, these cases must be addressed. Below, we define the concept of generalized solutions, which allows us to handle this issue without resorting to the theory of distributions.


\paragraph{Definition}

Let $\Omega$ be an open set with a $\mathcal{C}^1$ boundary in $\mathbb{R}^n$, and let $a$ and $S$ be continuous on $\mathbb{R}_+ \times \Omega$. We say that $f$ is a generalized solution of the problem

$$\frac{\partial f}{\partial t}(t,x)+v \cdot \nabla_x f(t,x)+ a(t,x)f(t,x)=S(t,x), \quad x \in \mathbb{R}^N, \quad t>0.\\$$

if and only if the function

$s \mapsto f(t+s,x+sv)$ is of class $\mathcal{C}^1$ almost everywhere in $(t,x) \in \mathbb{R}_+ \times \Omega$ and satisfies

$\frac{d}{ds}f(t+s,x+sv) +a(t+s,x+sv)f(t+s,x+sv) = S(t+s,x+sv) \quad (t+s,x+sv) \in \mathbb{R}_+^* \times \Omega$.


\paragraph{}



Let $\Omega$ be an open set with a $\mathcal{C}^1$ boundary in $\mathbb{R}^n$, and let $a$ and $S$ be continuous on $\mathbb{R}_+ \times \Omega$. Let $f_b^- \in \mathbf{L}_{loc}^{\infty}(\mathbb{R}_+ \times \partial \Omega^-)$ and $f^{in} \in \mathbf{L}_{loc}^{\infty}(\overline{\Omega})$.

In this case, the problem:

\[
\begin{cases}
	\frac{\partial f}{\partial t}(t,x)+v \cdot \nabla_x f(t,x)+ a(t,x)f(t,x)=S(t,x), \quad x \in \mathbb{R}^N, \quad t>0.\\
	\restriction{f}{\partial \Omega}= f_b^-\\
	\restriction{f}{t=0}= f^{in}(x)
\end{cases}
\]

has a unique generalized solution. The formula for this solution is the same as in the previous case, i.e., almost everywhere in $(t,x) \in \mathbb{R}_+ \times \Omega$, we have:

\[
\boxed{
	\begin{aligned}
		f(t,x) &= \mathbf{1}_{t \leq \tau_x} f^{in} (x-tv) \exp\left(-\int_0^t a(s,x+(s-t)v)ds\right) \\
		&\quad + \mathbf{1}_{t > \tau_x} f_b^- (t-\tau_x,x-\tau_x v) \exp\left(-\int_{t-\tau_x}^t a(s,x+(s-t)v)ds\right) \\
		&\quad + \int_{t-\tau_x}^t \exp\left(-\int_{s}^{t}a(\tau,x+(\tau-t)v)d\tau\right) S(s,x+(s-t)v)ds
\end{aligned}}
\]







\section{Linear Boltzmann Equation}

In this section, we will focus on the actual Boltzmann equation. This integro-differential equation is described in \ref{Boltzmann}.

\begin{equation} \label{Boltzmann}
	\frac{\partial f}{\partial t}(t,x,v)+v \cdot \nabla_x f(t,x,v) + a(t,x)f(t,x,v) = \int_{\mathbb{R}^n}k(t,x,v,w)f(t,x,w)d\mu(w)+Q(t,x,v).
\end{equation}

To simplify notations, this equation will be written as: 
\[ \frac{\partial f}{\partial t}+v \cdot \nabla_x f +af =   \mathcal{K} f +Q \]

Here, $v$ and $k$ are given continuous functions. The range of possibilities for the velocity $v$ is very diverse due to the variety of physical problems that can be considered. Therefore, the notation $d\mu(w)$ encompasses all these possibilities by incorporating a weight function. The major difference from the transport equation is that, this time, we consider exchange phenomena, and thus, $v$ cannot be seen as a simple parameter.

In this section, we cannot highlight the exact expression of the solutions because the method of characteristics yields another integral equation involving the solution.




\paragraph{Definition}

Let $Q$ be continuous on $]0,T[\times \Omega \times \mathbb{R}^n$. We say that $f$ is a generalized solution of the problem, continuous on $]0,T[\times \Omega \times \mathbb{R}^n$,


$$\frac{\partial f}{\partial t}+v \cdot \nabla_x f +af = \mathcal{K} f +Q\\$$

\paragraph{}
if and only if, for all $(t,x,v) \in ]0,T[\times \Omega \times \mathbb{R}^n$, the function 

\[ s \mapsto f(t+s,x+sv,v) \]

is of class $\mathcal{C}^1$ for all $x+sv \in \Omega$ and satisfies, for all $s \in \mathbb{R}$ such that $x +sv \in \Omega$,

$\frac{d}{ds}f(t+s,x+sv,v) +a(t+s,x+sv,v)f(t+s,x+sv,v) = (\mathcal{K}f+ Q)(t+s,x+sv,v).$

\subsection{Cauchy Problem}
\paragraph{}



Let $f^{in} \in \mathbf{C}_b(\mathbb{R}^n \times \mathbb{R}^n)$ and $Q \in \mathbf{C}_b(]0,T[ \times \mathbb{R}^n \times \mathbb{R}^n)$. Also, let $0 \leq a \in \mathbf{C}_b(]0,T[ \times \mathbb{R}_x^n \times \mathbb{R}_v^n)$ and $0 \leq k \in \mathbf{C}_b(]0,T[ \times \mathbb{R}_x^n \times \mathbb{R}_v^n \times \mathbb{R}_w^n)$. The notation $\mathbf{C}_b(X)$ denotes the set of bounded continuous functions on $X$ with values in $\mathbb{R}$.


We assume that:  \[ \underset{(t,x,v) \in ]0,T[ \times \mathbb{R}^n \times \mathbb{R}^n}{\sup} \int_{\mathbb{R}^n} k(t,x,v,w) \, d\mu(w) < +\infty \]

\paragraph{}

Then the problem is:

\[
\begin{cases}
	\frac{\partial f}{\partial t}+v \cdot \nabla_x f +af =   \mathcal{K} f +Q \quad (t,x,v) \in ]0,T[ \times \mathbb{R}^n \times \mathbb{R}^n\\
	\restriction{f}{t=0}= f^{in}
\end{cases}
\]
has a unique generalized solution $f \in \mathbf{C}_b(]0,T[ \times \mathbb{R}^n \times \mathbb{R}^n)$.

\subsection{Boundary Value Problem}

We define the set $\Gamma_-$ as $\Gamma_-= \{(x,v) \in \partial \Omega \times \mathbb{R}^n ~ | ~ v\cdot n_x<0\}$

Let $f^{in} \in \mathbf{C}_b(\overline{\Omega} \times \mathbb{R}^n)$, $f_b^- \in \mathbf{C}_b([0,T] \times \Gamma_-)$ such that $f_b^-(0,y,v)=f^{in}(y,v) \quad \forall (y,v) \in \Gamma_-$

Let $Q \in \mathbf{C}_b([0,T] \times \overline{\Omega} \times \mathbb{R}^n)$.

Let $0 \leq a \in \mathbf{C}_b(\mathbb{R}_+ \times \Omega_x \times \mathbb{R}_v^n)$ and $0 \leq k \in \mathbf{C}_b(\mathbb{R}_+ \times \Omega_x^n \times \mathbb{R}_v^n \times \mathbb{R}_w^n)$.

Then the problem

\[
\begin{cases}
	\frac{\partial f}{\partial t}+v \cdot \nabla_x f +af =   \mathcal{K} f +Q \quad (t,x,v) \in ]0,T[ \times \Omega \times \mathbb{R}^n\\
	\restriction{f}{t=0}= f^{in}\\
	\restriction{f}{\Gamma^-}=f_b^-
\end{cases}
\]

has a unique generalized solution $f \in \mathbf{C}_b([0,T] \times \Omega \times \mathbb{R}^n)$.

\section{The use of Monte Carlo methods to solve the linear Boltzmann equation}

The transport equation in an infinite medium with its corresponding deterministic collisional component can be expressed as:
\begin{equation}
	\partial _t u(x,t,\textbf{v}) + \textbf{v} \cdot \nabla u(x,t,\textbf{v}) + v\sigma_t (x,t,\textbf{v})u(x,t,\textbf{v})= v\sigma_s(x,t,\textbf{v})\int P (x,t,\textbf{v},\textbf{v}')u(x,t,\textbf{v}')d\textbf{v}' \label{ref11}
\end{equation}

Where 
\begin{equation*}
	\sigma_s (x,t,\textbf{v})= \int \sigma_s (x,t,\textbf{v},\textbf{v}')d\textbf{v}', \quad  P (x,t,\textbf{v},\textbf{v}')=
	\frac{\sigma_s (x,t,\textbf{v},\textbf{v}')}{\sigma_s (x,t,\textbf{v})}
\end{equation*}

The approach involves a series of variable changes. The initial step involves re-expressing the transport equation \ref{ref11} with respect to a characteristic $x + vt$. As a result, it transforms into:

\begin{equation}
	\partial _s u(x+\textbf{v}s,s,\textbf{v}) = -v\sigma_t (x+\textbf{v}s,s,\textbf{v})u(x+\textbf{v}s,s,\textbf{v}) + v\sigma_s(x+\textbf{v}s,s,\textbf{v})\int P (x+\textbf{v}s,s,\textbf{v},\textbf{v}')u(x+\textbf{v}s,s,\textbf{v}')d\textbf{v}'
\end{equation}

Let's multiply both sides of the equation by:
\begin{equation*}
	e^{\int _0^s v\sigma_t (x + \textbf{v}\alpha,\alpha, v) d\alpha}
\end{equation*}
Following that, we obtain
\begin{equation*}
	\partial _s [u(x+\textbf{v}s,s,\textbf{v})e^{\int _0^s v\sigma_t (x + \textbf{v}\alpha,\alpha, v) d\alpha}] = e^{\int _0^s v\sigma_t (x + \textbf{v}\alpha,\alpha, v) d\alpha} v\sigma_s(x+\textbf{v}s,s,\textbf{v})\int P (x+\textbf{v}s,s,\textbf{v},\textbf{v}')u(x+\textbf{v}s,s,\textbf{v}')d\textbf{v}' \label{ref12}
\end{equation*}
We get after integrating the equation \ref{ref12} between [0, t]:

\begin{multline}
	u(x+\textbf{v}t,t,\textbf{v}) = u_0(x, \textbf{v}) \exp\left(- \int_{0}^{t} v\sigma_t\left(x + \textbf{v} \alpha, \alpha, \textbf{v}\right) d\alpha\right) \\
	+ \int_{0}^{t} \int v\sigma_s\left(x + \textbf{v}s, s, \textbf{v}\right) u\left(x + \textbf{v}s, s, \textbf{v}'\right) e^{- \int_s^t v\sigma_t\left(x + \textbf{v} \alpha, \textbf{v}\right) d\alpha} P\left(x + \textbf{v} s, s, \textbf{v}, \textbf{v}'\right) d\textbf{v}'ds 
\end{multline}

We obtain:

\begin{multline}
	u(x,t,\textbf{v}) = u_0(x - \textbf{v}t, \textbf{v}) \exp\left(- \int_{0}^{t} v\sigma_t\left(x - \textbf{v}(t - \alpha), \alpha, \textbf{v}\right) d\alpha\right) \\
	+ \int_{0}^{t} \int v\sigma_s\left(x - \textbf{v}(t - s), s, \textbf{v}\right) u\left(x - \textbf{v}(t - s), s, \textbf{v}'\right) e^{- \int_s^t v\sigma_t\left(x - \textbf{v}(t - \alpha), \textbf{v}\right) d\alpha} P\left(x - \textbf{v}(t - s), s, \textbf{v}, \textbf{v}'\right) d\textbf{v}'ds \label{ref1}
\end{multline}
We also have:
\begin{multline}
	\exp\left(- \int_{0}^{t} v\sigma_t\left(x - \textbf{v}(t - \alpha), \alpha, \textbf{v}\right) d\alpha\right) = \exp\left(- \int_{0}^{t} v\sigma_t\left(x - \textbf{v} \alpha,t- \alpha, \textbf{v}\right) d\alpha\right) \\ = \int _t^\infty  v\sigma_t\left(x - \textbf{v} s,t- s, \textbf{v}\right)
	\exp\left(- \int_{0}^{s} v\sigma_t\left(x - \textbf{v} \alpha,t- \alpha, \textbf{v}\right) d\alpha\right) ds
\end{multline}
Then the integral representation of \ref{ref11} is provided by:

\begin{multline}
	u(x,t,\textbf{v}) =  \int _t^\infty  u_0(x - \textbf{v}t, \textbf{v}) v\sigma_t\left(x - \textbf{v} s,t- s, \textbf{v}\right)
	\exp\left(- \int_{0}^{s} v\sigma_t\left(x - \textbf{v} \alpha,t- \alpha, \textbf{v}\right) d\alpha\right) ds\\
	+ \int_{0}^{t} \int v\sigma_s\left(x - \textbf{v}(t - s), s, \textbf{v}\right) u\left(x - \textbf{v}(t - s), s, \textbf{v}'\right) e^{- \int_s^t v\sigma_t\left(x - \textbf{v}(t - \alpha), \textbf{v}\right) d\alpha} P\left(x - \textbf{v}(t - s), s, \textbf{v}, \textbf{v}'\right) d\textbf{v}'ds \label{ref13}
\end{multline}


\section{Semi analog MC scheme}

Developing a Monte Carlo scheme involves introducing a group of random variables along with their associated probability measure to express \ref{ref13} as an expectation. The selection of this set of random variables is not unique, resulting in various Monte Carlo schemes with distinct properties.

\section{Numerical tests/results}

à voir

\section{Conclusion}


	
	
\end{document}